\begin{table}[H]
\centering
\sffamily
\small
\renewcommand{\arraystretch}{1.15}
\begin{threeparttable}
\caption{D\'ecomposition internationale de la croissance de la productivit\'e du travail, 2001--2019 (\%)}
\label{tab:note_international}
\begin{tabular}{l *{4}{r}}
\toprule
& $\Delta \ln(Y/L)$ & PTF intra & R\'ealloc. & Capital \\
\midrule
Lettonie & 5.90 & -1.79 & 1.00 & 6.68 \\
Lituanie & 4.13 & -0.44 & -0.19 & 4.76 \\
Estonie & 3.58 & -1.30 & 0.64 & 4.24 \\
Tch\'equie & 2.76 & 2.37 & -0.11 & 0.50 \\
Slov\'enie & 1.86 & 0.59 & -0.67 & 1.94 \\
\'Etats-Unis & 1.60 & 0.90 & 0.04 & 0.66 \\
Su\`ede & 1.36 & 0.37 & -0.17 & 1.16 \\
Danemark & 1.17 & 0.32 & -0.13 & 0.99 \\
Allemagne & 1.14 & 0.92 & 0.01 & 0.21 \\
Autriche & 1.09 & 1.02 & -0.03 & 0.10 \\
Royaume-Uni & 1.03 & 0.62 & -0.18 & 0.60 \\
Belgique & 1.00 & 0.54 & -0.08 & 0.53 \\
Japon & 0.95 & 0.55 & -0.01 & 0.42 \\
Pays-Bas & 0.94 & 0.19 & -0.02 & 0.76 \\
Finlande & 0.90 & 1.12 & -0.32 & 0.09 \\
France & 0.83 & 0.61 & -0.06 & 0.28 \\
Espagne & 0.80 & -0.51 & 0.01 & 1.30 \\
\textbf{Canada} & \textbf{0.52} & \textbf{0.12} & \textbf{-0.26} & \textbf{0.65} \\
Italie & 0.19 & -0.90 & -0.10 & 1.19 \\
\bottomrule
\end{tabular}
\begin{tablenotes}\footnotesize
\item \textit{Note}\,: Croissance annuelle moyenne en points de pourcentage, 2001--2019. Classification\,: 15 secteurs NACE (O et P exclus). La croissance agr\'eg\'ee de la productivit\'e du travail est calcul\'ee comme la croissance de la VA moins celle du travail, chacune agr\'eg\'ee par indices de T\"ornqvist sur les m\^emes 15 secteurs. La PTF intra-industries et la r\'eallocation sont amplifi\'ees par $1/(1-\bar\alpha)$; la composante capital est le r\'esidu. Source\,: EU-KLEMS 2025, module Growth Accounts Basic.
\end{tablenotes}
\end{threeparttable}
\end{table}